% Decision process, separate from the workflow construction layer.
\subsection{A1：先分开短段策略、方向与 DP 的作用}
A 读取同一 240 条开发记录的已核验对照后，将当时方案从 R0 更新为 S0，触发的是 107 条覆盖改善与全部保护通过。点数、长度两项各自具有作用，但单独保留任一项都会损失组合恢复的一部分覆盖，因此没有为了更少参数把必要调整删去。

统一方向 $60^\circ$ 的 3 条几何退化被保留为真实负结果。A 没有反复试到统一方向胜出，而是核对事先登记的原始时间分组机制：时间规则组有局部收益，异常/不可计算组回到基线。这只支持下一项真实检查，不提前证明未见记录安全。

\subsection{A2：有用父项不等于立即替换}
G0 相对 R0 的保护通过，41 条具有几何改善；与 S0 比较时，它会丢失 6,704 个新增覆盖点，S0 又不能保留 G0 的全部几何收益。按预先批准的保守规则，当时方案仍是 S0。两部分作用机制不同且已独立核验，才允许进入 \texttt{R0/S0/G0/S0\_G0} 四种完整对照。

\subsection{A3：接受实际组合，拒绝无贡献的附加项}
实际组合 \texttt{S0\_G0} 同时通过对固定 R0 和当时 S0 的全部保护，相对 S0 在 41 条上改善共同几何误差，开发中的保留方案才更新为组合。移去 G0 会损失这些已接受性质；移去 S0 会损失 6,704 个覆盖身份。这个决定来自父项与移除比较，不来自预设 S0 必胜或“组件越多越好”。

S 与 DP 的两个组合也真实执行：DP 2 在相同上游多存 2,010 点，整链相对 S0 有 9 条几何退化；DP 10 在 178 条超过发布预算。A 保留这些负结果，未把覆盖增长解释成 P 改进，最终开发短名单仍使用 DP 5。

剩余机会复盘认为，有依据的父项、主要交互与必要移除已经完成。未采用第二类删除保护，因为没有可靠的位移/物理真值阈值；未重跑完整四模式矩阵，因为历史材料尚不支持增加记录级模型或记忆依赖。实际只新增 1 个轻量单项与 3 个组合，没有为了耗尽配额补空实验。开发中的路径为 R0 $\rightarrow$ S0 $\rightarrow$ S0 $\rightarrow$ \texttt{S0\_G0}；短名单进入下一阶段，尚不能代替最终确认。

这些都是在用户事先批准规则内的系统决定。独立 C 已对本轮开发产物与裁决闭合检查，但内部核验不是新的用户批准，也不是 Evidence Master 的互动证据 Lock。没有可用的用户逐次判断原话，因此不补写“用户选中了组合”。
